\chapter{Impacto y responsabilidad}
\label{chap:impacto}

Aplicar inteligencia artificial a la información jurídica no es una decisión técnica neutra: afecta a cómo la ciudadanía accede a sus derechos, moviliza recursos económicos y computacionales, y activa responsabilidades legales y éticas. El caso es además especialmente sensible porque la tecnología —los modelos de lenguaje y la arquitectura RAG— es reciente y su marco regulador todavía está cristalizando. Este capítulo analiza esos impactos con la misma honestidad que ha guiado el resto del trabajo, sin presentar el sistema como más benigno de lo que es. Conviene recordar la premisa que condiciona todo el análisis: el sistema es una \emph{herramienta informativa de apoyo a la consulta documental}, no un servicio de asesoramiento jurídico ni un instrumento de decisión con efectos legales.

\section{Impacto social}
\label{sec:impacto-social}

El principal impacto social positivo es de \textbf{acceso}. La información jurídica está publicada, pero no por ello es accesible: su estructura, su lenguaje técnico y su volumen levantan una barrera real para quien no tiene formación jurídica. Un sistema que acepta preguntas en lenguaje cotidiano, responde de forma comprensible y remite a la fuente oficial acerca la norma a la ciudadanía, en la línea de las iniciativas que ven en la IA una vía para mejorar el acceso a la justicia \cite{oecdAccessJusticeAI}. El diseño refuerza este objetivo con dos rasgos: un lenguaje llano y la trazabilidad de cada respuesta a su artículo de origen, que permite a la persona verificar y profundizar.

Ese mismo poder encierra los riesgos. El primero es la \textbf{falsa sensación de certeza}: una respuesta fluida y bien citada puede parecer definitiva a quien no tiene criterio para detectar un error sutil, un problema documentado en las herramientas jurídicas de IA \cite{magesh2025hallucinationfree}. El segundo es el \textbf{uso indebido}: emplear el sistema como si fuera un dictamen, o para decidir cuestiones con consecuencias jurídicas. El diseño mitiga ambos —aviso constante sobre su carácter informativo, abstención cuando falta evidencia, negativa a presentarse como asesoramiento—, pero no los elimina; la responsabilidad última de la interpretación sigue siendo humana. A ellos se añaden los riesgos generales de los modelos de lenguaje —sesgos y representaciones desiguales— señalados por la literatura crítica \cite{bender2021stochastic}.

En cuanto a la \textbf{igualdad}, el sistema puede reducir la desigualdad de acceso a la información normativa —especialmente para quien no puede costear asesoramiento—, aunque no debe ignorarse que su uso presupone acceso tecnológico, de modo que, mal desplegado, podría reproducir la brecha digital en lugar de cerrarla. Finalmente, un aspecto de \textbf{privacidad} con dimensión social: las consultas jurídicas revelan a menudo situaciones personales delicadas —un despido, una deuda, un expediente de extranjería—; mantener el procesamiento en local, como hace este sistema, evita que esas preguntas salgan hacia servicios de terceros.

\section{Impacto económico}
\label{sec:impacto-economico}

La decisión de construir el sistema íntegramente sobre \textbf{software y modelos de código abierto ejecutados en local} \cite{touvron2023llama2} tiene consecuencias económicas directas. La más evidente es la ausencia de \textbf{coste por consulta}: frente a las \emph{APIs} comerciales, aquí no se paga por cada pregunta, ni existe dependencia de un proveedor externo que pueda cambiar precios, condiciones o disponibilidad. El coste se concentra en el hardware del servidor y en el tiempo de desarrollo, no en una factura recurrente.

En \textbf{mantenimiento}, el diseño reproducible y el corpus ampliable por configuración abaratan la evolución del sistema: incorporar nuevas normas es un proceso administrativo, no un rediseño. La \textbf{dependencia de servicios externos} es, por tanto, prácticamente nula, lo que además protege frente al bloqueo por proveedor (\emph{lock-in}). En cuanto a \textbf{escalabilidad}, el índice exacto es sobrado para la escala del corpus; crecer en órdenes de magnitud exigiría técnicas de búsqueda aproximada, viables pero hoy innecesarias. El límite económico real es la \textbf{latencia en CPU}: la inferencia local sin GPU es lenta, lo que condiciona la viabilidad de un servicio masivo en tiempo real y obligaría, en producción, a invertir en aceleración o a asumir tiempos de respuesta altos. Como prototipo demostrable y reproducible, sin embargo, es plenamente viable en hardware modesto.

\section{Impacto medioambiental}
\label{sec:impacto-medioambiental}

Conviene ser honesto y preciso, evitando tanto la alarma como el «lavado verde». El mayor coste ambiental de los modelos de lenguaje está en su \textbf{entrenamiento}, que consume cantidades notables de energía \cite{strubell2019energy}. Este trabajo \textbf{no entrena ningún modelo}: reutiliza modelos abiertos ya preentrenados, con lo que evita por completo esa huella, que es la dominante. El coste ambiental del proyecto se limita, pues, a la \textbf{inferencia} y a operaciones puntuales de indexación.

Esa inferencia corre en \textbf{CPU} y se apoya en modelos \textbf{cuantizados} \cite{frantar2023gptq}, que reducen de forma drástica la memoria y el cómputo necesarios. El resultado es un consumo por consulta modesto, si bien con un matiz honesto: la CPU es más lenta que la GPU, de modo que, aunque su potencia instantánea sea menor, el tiempo por consulta es mayor. El \textbf{almacenamiento} es despreciable —el corpus es de 92 normas y el índice, de decenas de miles de vectores, no de millones—, y la decisión de usar un índice exacto en lugar de uno aproximado no penaliza a esta escala. En conjunto, reutilizar en lugar de entrenar, cuantizar y ejecutar en local es la dirección responsable; pero sería deshonesto negar que toda IA tiene una huella y que un despliegue a gran escala la multiplicaría.

\section{Responsabilidad ética y profesional}
\label{sec:responsabilidad}

El marco regulador es el punto de partida. El \textbf{Reglamento (UE) 2024/1689}, la Ley de Inteligencia Artificial de la Unión, adopta un enfoque basado en el riesgo \cite{euAiAct2024}. Un asistente meramente \emph{informativo}, que no toma decisiones con efectos jurídicos ni sustituye a un profesional, no encaja en la categoría de «alto riesgo», pero sí queda sujeto a obligaciones de \textbf{transparencia}: la persona debe saber que interactúa con un sistema automático y cuáles son sus límites. El tratamiento de datos personales, a su vez, se rige por el \textbf{Reglamento General de Protección de Datos} \cite{gdpr2016}; como las consultas pueden contener datos personales, la ejecución local —sin envío a terceros— es una salvaguarda de partida. A ello se suman los principios de la Carta Ética Europea sobre el uso de la IA en los sistemas judiciales \cite{cepejEthicalCharter}.

Sobre esa base, el diseño asume varias responsabilidades de forma explícita. Frente al \textbf{riesgo de alucinación} —bien documentado en el dominio jurídico \cite{dahl2024legalhallucinations,magesh2025hallucinationfree} y que la propia arquitectura RAG reduce pero no elimina \cite{kim2025legalragfail}—, el sistema opone trazabilidad, citas verificables, abstención ante evidencia insuficiente y validación estricta de la salida. La \textbf{transparencia} se materializa en que cada afirmación remite a su fuente y en que se documenta qué versión del corpus se utilizó. En cuanto a la \textbf{propiedad intelectual}, los textos proceden de los datos abiertos del BOE, los modelos son de pesos abiertos con sus respectivas licencias y el código es de código abierto. Y la \textbf{supervisión humana} no es opcional: el sistema es un apoyo, no un sustituto, y su propia salida lo recuerda mediante un aviso que subraya que los textos consolidados del BOE tienen valor informativo y no oficial \cite{boeLegislacionActualizada}. Esta combinación —límites declarados, fuentes verificables y última palabra humana— es la forma concreta en que el trabajo traduce la responsabilidad profesional en decisiones de diseño.
